\documentclass[11pt,a4paper,oneside]{book}
\usepackage[utf8]{inputenc}
\usepackage[T1]{fontenc}
\usepackage[brazil]{babel}
\usepackage{lmodern,microtype}
\usepackage{geometry}
\geometry{margin=2.6cm}
\usepackage{amsmath,amssymb,mathtools}
\usepackage{siunitx,booktabs,hyperref,listings,xcolor}
\hypersetup{colorlinks=true,linkcolor=blue,urlcolor=blue}
\title{Engenharia Aplicada à Perfusão\\\large Modelagem matemática, métodos numéricos, CFD e geometria}
\author{Gabriel Passos de Jesus\\\small PerfusionLab}
\date{Versão de trabalho — novembro de 2026}
\begin{document}
\frontmatter
\maketitle
\chapter*{Escopo e responsabilidade}
Esta apostila apresenta fundamentos matemáticos e computacionais para investigação de fenômenos relacionados à perfusão e à circulação extracorpórea. Os modelos são abstrações educacionais. Não representam dispositivos comerciais, anatomia individual, sangue em toda a sua complexidade ou condições clínicas validadas. A implementação não deve orientar cuidados ou parâmetros de pacientes.
\tableofcontents
\mainmatter
\chapter{Formulação de problemas}
\section{Sistema, domínio e observáveis}
Todo modelo começa pela definição do domínio físico, fronteiras, variáveis de estado, entradas, parâmetros e observáveis. É necessário declarar escalas espaciais e temporais, hipóteses constitutivas e fontes de incerteza. Uma saída numérica sem unidades e contexto não é um resultado científico interpretável.

\section{Análise dimensional}
As dimensões fundamentais são massa $M$, comprimento $L$ e tempo $T$. Pressão possui dimensão $ML^{-1}T^{-2}$, vazão volumétrica $L^3T^{-1}$, viscosidade dinâmica $ML^{-1}T^{-1}$ e densidade $ML^{-3}$. A análise dimensional detecta inconsistências e orienta grupos adimensionais.

\chapter{Conservação}
\section{Massa}
A forma integral da conservação de massa é
\[
\frac{d}{dt}\int_{\Omega}\rho\,dV+
\int_{\partial\Omega}\rho\mathbf{u}\cdot\mathbf{n}\,dA=0.
\]
Para fluido incompressível de densidade constante, obtém-se $\nabla\cdot\mathbf{u}=0$.

\section{Quantidade de movimento}
A equação de Navier--Stokes para fluido newtoniano incompressível é
\[
\rho\left(\frac{\partial\mathbf{u}}{\partial t}
+\mathbf{u}\cdot\nabla\mathbf{u}\right)
=-\nabla p+\mu\nabla^2\mathbf{u}+\rho\mathbf{f}.
\]
A escolha de termos, constituição e condições de contorno depende do problema. Esta equação não torna automaticamente um modelo apropriado ao sangue ou a um dispositivo específico.

\chapter{Modelos concentrados: EDOs}
Um exemplo de dinâmica de primeira ordem para uma variável de fluxo $Q$ é
\[
\tau_Q\frac{dQ}{dt}=Q_{cmd}-Q.
\]
Para entrada constante, a solução analítica é
\[
Q(t)=Q_{cmd}+[Q(0)-Q_{cmd}]e^{-t/\tau_Q}.
\]
A constante de tempo controla a escala de relaxamento. O modelo é útil para testar integradores e não pretende reproduzir toda a dinâmica cardiovascular.

\section{Runge--Kutta de quarta ordem}
Para $y'=f(t,y)$, o método RK4 usa quatro avaliações por passo:
\[
\begin{aligned}
k_1&=f(t_n,y_n),\\
k_2&=f(t_n+h/2,y_n+hk_1/2),\\
k_3&=f(t_n+h/2,y_n+hk_2/2),\\
k_4&=f(t_n+h,y_n+hk_3),\\
y_{n+1}&=y_n+\frac{h}{6}(k_1+2k_2+2k_3+k_4).
\end{aligned}
\]
A ordem formal não elimina a necessidade de estudo de convergência e controle de erro.

\chapter{Modelos distribuídos: EDPs}
\section{Advecção--difusão--reação}
Uma variável escalar transportada pode ser descrita por
\[
\frac{\partial C}{\partial t}
+u\frac{\partial C}{\partial x}
=D\frac{\partial^2 C}{\partial x^2}-k(C-C_{eq}).
\]
O termo advectivo transporta a variável com o campo de velocidade; o difusivo suaviza gradientes; o termo reacional representa uma aproximação fenomenológica de relaxamento. A variável $C$ não é, por padrão, concentração fisiológica medida.

\section{Condições de contorno}
Condição de Dirichlet especifica o valor da variável na fronteira. Neumann especifica sua derivada normal. A escolha precisa representar o problema matemático e ser compatível com o sentido do escoamento. Condições iniciais e de contorno devem ser testadas isoladamente.

\chapter{Métodos numéricos}
\section{Diferenças finitas e volumes finitos}
Diferenças finitas aproximam derivadas em nós. Volumes finitos integram leis de conservação em células e contabilizam fluxos pelas faces. A segunda abordagem é especialmente útil quando o balanço conservativo é central.

\section{Advecção upwind e difusão central}
O esquema upwind escolhe a diferença na direção do escoamento e tende a ser robusto, embora possa introduzir difusão numérica. A difusão pode ser aproximada por diferenças centrais em malha regular. A ordem, estabilidade e erro devem ser documentados.

\section{SSP-RK2}
Uma forma de segunda ordem é
\[
y^{(1)}=y^n+\Delta t F(y^n),\qquad
y^{n+1}=\frac12 y^n+\frac12\left[y^{(1)}+\Delta t F(y^{(1)})\right].
\]
A propriedade SSP depende da estabilidade da etapa Euler e da restrição de passo correspondente.

\chapter{Estabilidade e convergência}
Números adimensionais como Reynolds, Peclet e Courant ajudam a caracterizar regimes e restrições numéricas:
\[
Re=\frac{\rho U L}{\mu},\qquad
Pe=\frac{UL}{D},\qquad
Co=\frac{U\Delta t}{\Delta x}.
\]
Não existe um único limite universal aplicável a todos os esquemas. O limite deve ser derivado para a discretização efetivamente implementada.

Verificações recomendadas: solução analítica, solução manufaturada, conservação, refinamento espacial, refinamento temporal, ordem observada, sensibilidade paramétrica e testes de regressão.

\chapter{CFD: primeiro problema de referência}
Para escoamento laminar, incompressível, newtoniano, estacionário e totalmente desenvolvido em tubo reto rígido, a solução de Poiseuille é
\[
u(r)=\frac{G}{4\mu}(R^2-r^2).
\]
O PerfusionLab inicia com a equação transiente radial para velocidade axial, discretizada por volumes finitos. Esse problema reduzido permite verificar simetria, condição de não deslizamento, resposta ao gradiente de pressão e aproximação da solução estacionária. Ele não é um solver Navier--Stokes 3D nem representa uma cânula real.

\chapter{Geometria computacional e modelagem 3D}
A geometria define o domínio sobre o qual as equações serão resolvidas. Parâmetros como diâmetro, comprimento, curvatura, espessura e conexões devem ter origem documentada. Um modelo visual não é automaticamente uma geometria de engenharia, uma malha de qualidade ou um dispositivo fabricável.

Fluxo de trabalho sugerido:
\begin{enumerate}
\item Definir finalidade e escala da peça.
\item Parametrizar dimensões e tolerâncias conceituais.
\item Gerar sólido CAD e verificar unidades.
\item Exportar formato de intercâmbio e conferir integridade.
\item Gerar malha, avaliar qualidade e independência de malha.
\item Definir propriedades, condições de contorno e solver.
\item Comparar com solução analítica ou benchmark.
\end{enumerate}

\chapter{Acoplamento ODE--PDE--CFD}
O acoplamento deve ocorrer por contratos explícitos de variáveis e unidades. Um fluxo agregado calculado por ODE pode ser convertido em velocidade média apenas quando a área efetiva é definida e suas hipóteses são claras. O campo CFD pode futuramente fornecer relações de perda de carga ou perfis para modelos reduzidos, mas a retroalimentação exige verificação e validação independentes.

\chapter{Verificação, validação e incerteza}
\textbf{Verificação} pergunta se as equações foram resolvidas corretamente. \textbf{Validação} pergunta se as equações representam adequadamente o fenômeno de interesse para um uso definido. Calibração não substitui validação. Documentar versão, malha, passo temporal, tolerâncias, dados, parâmetros, ambiente e limitações é parte do resultado.

\chapter{Roteiro prático do repositório}
\begin{enumerate}
\item Executar testes Java/Maven do núcleo ODE/PDE.
\item Executar os testes Python em \texttt{cfd/}.
\item Inspecionar a geometria OpenSCAD em \texttt{geometry/}.
\item Reproduzir os benchmarks e guardar parâmetros e logs.
\item Registrar hipóteses e evidências em documentação versionada.
\end{enumerate}
\backmatter
\chapter*{Referências para aprofundamento}
\begin{thebibliography}{9}
\bibitem{white} White, F. M. \textit{Fluid Mechanics}. McGraw-Hill.
\bibitem{versteeg} Versteeg, H. K.; Malalasekera, W. \textit{An Introduction to Computational Fluid Dynamics: The Finite Volume Method}. Pearson.
\bibitem{ferziger} Ferziger, J. H.; Perić, M.; Street, R. L. \textit{Computational Methods for Fluid Dynamics}. Springer.
\bibitem{leveque} LeVeque, R. J. \textit{Finite Volume Methods for Hyperbolic Problems}. Cambridge University Press.
\bibitem{pope} Pope, S. B. \textit{Turbulent Flows}. Cambridge University Press.
\end{thebibliography}
\end{document}
